
\section{Discussion}
\label{sec:discussion}

\subsection{Explaining the Main Findings}

\textbf{Why does CTQRS exceed the threshold by a large margin?}
Qwen2.5-7B 3-shot achieves median CTQRS=88.89\%, surpassing
the 75\% GPT-4o reference by 13.89 percentage points with a
large effect ($\delta=0.908$). CTQRS measures structural
completeness through rule-based sub-checks (presence of S2R,
EB, AB, environment information, interface elements). The
3-shot Alpaca-LoRA prompt explicitly instructs the model to
produce four sections, and Qwen reliably generates all required
structural components. This structural conformance is consistent
with E02's observation that even fine-tuned models achieve
high CTQRS scores once they learn the target
template~\cite{acharya2025bugquality}.

\textbf{Why does SBERT fall substantially below the threshold?}
Our length analysis reveals that Qwen outputs are on average
1.82$\times$ longer than the Bugzilla ground-truth reports
(245.5 words vs.\ 162.0 words). Qwen produces verbose,
markdown-formatted responses with section headers (e.g.,
\texttt{\#\#\#\# Steps to Reproduce}) and detailed bullet
points, whereas Bugzilla ground-truth reports are written in
plain, concise technical prose without headers. The
\texttt{paraphrase-MiniLM-L6-v2} encoder encodes the entire
passage into a fixed-size vector, making cosine similarity
sensitive to differences in length and formatting style
beyond semantic content. E02 itself notes that ``differences
in the length of the bug report generated can lead to unfair
comparisons in n-gram overlap
metrics''~\cite{acharya2025bugquality}; our results show
this effect extends to SBERT as well.

To confirm that the SBERT gap is caused by style rather than
semantic content, we conducted an ablation study adding a
length constraint to the system prompt (``keep total output
under 150 words''). While this shortened Qwen outputs to a
mean of 109.1 words (0.70$\times$ GT), SBERT further
decreased from 0.632 to 0.609. This counter-intuitive
result confirms that the problem is \emph{style mismatch},
not verbosity alone: outputs shorter than GT are penalized
in the opposite direction. ROUGE-1 improved from 0.491 to
0.532 with the length constraint, consistent with higher
unigram overlap when outputs are more concise.

\textbf{Why does ROUGE-1 achieve only a small effect?}
Median ROUGE-1=0.491 exceeds the 0.47 threshold but with a
small effect size ($\delta=0.191$). The margin above threshold
is narrow (0.021 F1 points), suggesting lexical overlap is
adequate but not strongly superior to the GPT-4o baseline.
This is consistent with E02's ROUGE results for fine-tuned
models on the S2R section (ROUGE=0.52), where the detailed
and lengthy S2R content also suppressed overlap
scores~\cite{acharya2025bugquality}.

\subsection{Comparison with Prior Work}

Our CTQRS result (88.89\%) substantially exceeds those
reported for fine-tuned models: Qwen2.5-7B fine-tuned
achieves 77\%~\cite{acharya2025bugquality}, AgentReport
with QLoRA achieves 80.5\%~\cite{choi2025agentreport},
and PPO reinforcement learning yields
76\%~\cite{yang2025ctqrsrl}. Three factors explain this:
(1) CTQRS measures structural completeness, which few-shot
prompting achieves readily when the prompt template
explicitly names all required sections; (2) our corrected
scorer (\texttt{max\_possible=18}) differs from the
original GindeLab scorer (\texttt{max\_possible=16}),
making direct numeric comparison with prior work
approximate; (3) fine-tuned models may produce shorter,
stylistically closer outputs that trade CTQRS for better
SBERT.

Our SBERT (0.632) is lower than the GPT-4o reference
(0.82) and likely lower than fine-tuned models, which
are trained to reproduce the exact style of Bugzilla
ground-truth reports. This confirms that few-shot
prompting, while effective for structural quality, cannot
replicate the writing style of domain-specific fine-tuned
models without additional training.

\subsection{Implications}

\textbf{For budget-constrained teams.}
Qwen2.5-7B 3-shot prompting --- requiring no GPU training,
no labeled data, and zero API cost when run on free cloud
compute --- produces structurally complete bug reports
(CTQRS=88.89\%) and adequate lexical fidelity
(ROUGE-1=0.491) that both exceed published GPT-4o few-shot
benchmarks. Teams lacking the resources to fine-tune a
model or subscribe to a proprietary API can adopt this
zero-training approach without sacrificing structural
quality. The only limitation is output verbosity, which
affects SBERT but not the structural completeness
perceived by developers reviewing the report.

\textbf{For researchers choosing evaluation metrics.}
SBERT similarity is sensitive to stylistic differences
between generated and reference texts beyond semantic
content. When the generation model and the reference corpus
have systematically different writing styles, SBERT scores
reflect style mismatch rather than content quality.
Researchers evaluating LLM-generated bug reports should
complement SBERT with style-agnostic metrics such as CTQRS
or human evaluation when the target style is not enforced
by the prompt.

\textbf{For practitioners choosing between few-shot and fine-tuning.}
Our findings support a practical decision rule: use
few-shot prompting when structural completeness (CTQRS) is
the primary requirement and compute budget is limited;
invest in fine-tuning when style alignment with a specific
reference corpus (SBERT) is required. Prior
work~\cite{acharya2025bugquality, choi2025agentreport,
yang2025ctqrsrl} shows that fine-tuned models achieve
higher SBERT by learning to replicate the exact writing
style of Bugzilla reports. The two approaches are
complementary: few-shot prompting provides a useful
starting point that fine-tuning can then refine for style.
